% ============================================================
% Sección 3
% ============================================================
\section{3 Consideraciones en la solución de una EDO}

\begin{tcolorbox}[enunciado]
\textbf{Clases:}
\begin{itemize}[leftmargin=*]
    \item Solución EDO con y sin unicidad 24Ago24 MCA4
    \item Sol Equilibrio, Curva Integral, Campo Direccional 26 Ago 2026 MCA4
\end{itemize}

\textbf{Software:}
\begin{itemize}[leftmargin=*]
    \item Wolfram Cloud: \texttt{EDO Logistica MCA4.nb}
    \item Github: \texttt{1.1\_Campo Direccional\_y\_Curva Integral}
\end{itemize}
\end{tcolorbox}

\vspace{0.3cm}
\textbf{Responde:}

\begin{enumerate}[label=\arabic*., leftmargin=*]
    \item Qué se entiende por curva integral (solución)
    
    \item Usa la gráfica (código) al final de \texttt{EDO Logistica\_MCA4.nb} o en el segundos 40 de la clase del 24 de Agosto. La gráfica muestra varias soluciones para la EDO Logística con distintas condiciones iniciales $P(t_{0})=P_{0}$.
    \begin{itemize}[leftmargin=*]
        \item ¿Cuál es la capacidad de carga en el código?
        \item ¿Qué tiempo $t_{0}$ iniciales se utilizaron?
        \item !!! Para el caso $y[0]=-1$ qué error tiene la solución que graficó el programa
    \end{itemize}
    \begin{enumerate}[label=(\alph*), leftmargin=*]
        \item Di en que caso se daría cada una de las rectas (punteadas) que ahí se muestran.
    \end{enumerate}
    
    \begin{figure}[H]
        \centering
        \includegraphics[width=0.8\textwidth]{figura1.png}
        \caption{Curvas Integrales de la Logística}
    \end{figure}

    \item LA EDO
    \[ \frac{dy}{dx}=\frac{4-2x}{3y^{2}-5} \]
    tiene por solución
    \[ y^{3}-5y+x^{2}-4x=c \]
    con $c \in \mathbb{R}$ una constante. Si graficamos esta expresión para algunas $c$, obtenemos curvas como se muestran en la figura 2.
    \begin{figure}[H]
        \centering
        \includegraphics[width=0.6\textwidth]{figura2.png}
        \caption{Gráficas de la Solución $y^3-5y+x^2-4x=c$}
    \end{figure}
    
    Trascribe los código \texttt{CURVAS\_EC\_SOL\_EDO\_DY\_4-2X\_DIV\_3Y\textasciicircum 2-5.PY}, \texttt{CAMPO\_DIRECCIONAL CURVAS 2X\_DIV\_3Y\textasciicircum 2-5.PY} y contesta lo siguiente:
    \begin{enumerate}[label=(\alph*), leftmargin=*]
        \item La curva dibujada para alguna $c$ es realmente una curva integral? justifica
        \item Para $c=-6$, gráficamente cuáles son las curvas integrales.
        \item Da condiciones iniciales que garanticen obtener cada una de las posibles soluciones del inciso anterior.
    \end{enumerate}
\end{enumerate}